\documentclass[CJK]{beamer}
%\usepackage{CJKutf8}
\usepackage{beamerthemesplit}
\usepackage{}
%\usetheme{Frankfurt}
%\usetheme{CambridgeUS}
%\usecolortheme{beaver}
%\usetheme{AnnArbor}
\usetheme{Dresden}
%\usebeamercolor{beetle}
\usepackage{xeCJK}
\setCJKmainfont{AR PL KaitiM GB}
%\useoutertheme{miniframes}
\usepackage{amsmath}
\usepackage{graphicx}
\usepackage{float} 
\usepackage{subfigure}
\usepackage{amssymb}
\usepackage{graphicx}
\usepackage{eufrak}
\usepackage{color}
\usepackage{array}
\usepackage{slashed}
\usepackage{simplewick}
\usepackage{tikz}
\usepackage{tcolorbox}
\usepackage[T1]{fontenc}
\graphicspath{{../figures/}}

%%figures
\def\lfig#1#2{\includegraphics[width=#1 in]{#2}}
\def\addfig#1#2{\begin{center}\includegraphics[width=#1 in]{#2}\end{center}}
\def\wulian{\includegraphics[width=0.18in]{emoji_wulian.jpg}}
\def\bigwulian{\includegraphics[width=0.35in]{emoji_wulian.jpg}}
\def\bye{\includegraphics[width=0.18in]{emoji_bye.jpg}}
\def\bigbye{\includegraphics[width=0.35in]{emoji_bye.jpg}}
\def\huaixiao{\includegraphics[width=0.18in]{emoji_huaixiao.jpg}}
\def\bighuaixiao{\includegraphics[width=0.35in]{emoji_huaixiao.jpg}}
\def\jianxiao{\includegraphics[width=0.18in]{emoji_jianxiao.jpg}}
\def\bigjianxiao{\includegraphics[width=0.35in]{emoji_jianxiao.jpg}}
%% colors
\def\blacktext#1{{\color{black}#1}}
\def\bluetext#1{{\color{blue}#1}}
\def\redtext#1{{\color{red}#1}}
\def\darkbluetext#1{{\color[rgb]{0,0.2,0.6}#1}}
\def\skybluetext#1{{\color[rgb]{0.2,0.7,1.}#1}}
\def\cyantext#1{{\color[rgb]{0.,0.5,0.5}#1}}
\def\greentext#1{{\color[rgb]{0,0.7,0.1}#1}}
\def\darkgray{\color[rgb]{0.2,0.2,0.2}}
\def\lightgray{\color[rgb]{0.6,0.6,0.6}}
\def\gray{\color[rgb]{0.4,0.4,0.4}}
\def\blue{\color{blue}}
\def\red{\color{red}}
\def\green{\color{green}}
\def\darkgreen{\color[rgb]{0,0.4,0.1}}
\def\darkblue{\color[rgb]{0,0.2,0.6}}
\def\skyblue{\color[rgb]{0.2,0.7,1.}}
%%control
\def\be{\begin{equation}}
\def\ee{\nonumber\end{equation}}
\def\bea{\begin{eqnarray*}}
\def\eea{\nonumber\end{eqnarray*}}
\def\bch{}
\def\ech{}
\def\bitem{\begin{itemize}}
\def\eitem{\end{itemize}}
\def\bcenter{\begin{center}}
\def\ecenter{\end{center}}
\def\bex{\begin{minipage}{0.2\textwidth}\includegraphics[width=0.6in]{jugelizi.png}\end{minipage}\begin{minipage}{0.76\textwidth}}
\def\eex{\end{minipage}}
\def\chtitle#1{\frametitle{\bch#1\ech}}
\def\bmat#1{\left(\begin{array}{#1}}
\def\emat{\end{array}\right)}
\def\bcase#1{\left\{\begin{array}{#1}}
\def\ecase{\end{array}\right.}
\def\bmini#1{\begin{minipage}{#1\textwidth}}
\def\emini{\end{minipage}}
\def\tbox#1{\begin{tcolorbox}#1\end{tcolorbox}}
\def\pfrac#1#2#3{\left(\frac{\partial #1}{\partial #2}\right)_{#3}}
%%symbols
\def\bropt{\,(\ \ \ )}
\def\sone{$\star$}
\def\stwo{$\star\star$}
\def\sthree{$\star\star\star$}
\def\sfour{$\star\star\star\star$}
\def\sfive{$\star\star\star\star\star$}
\def\rint{{\int_\leftrightarrow}}
\def\roint{{\oint_\leftrightarrow}}
\def\stdHf{{\textit{\r H}_f}}
\def\deltaH{{\Delta \textit{\r H}}}
\def\ii{{\dot{\imath}}}
\def\skipline{{\vskip0.1in}}
\def\skiplines{{\vskip0.2in}}
\def\lagr{{\mathcal{L}}}
\def\hamil{{\mathcal{H}}}
\def\vecv{{\mathbf{v}}}
\def\vecx{{\mathbf{x}}}
\def\vecy{{\mathbf{y}}}
\def\veck{{\mathbf{k}}}
\def\vecp{{\mathbf{p}}}
\def\vecn{{\mathbf{n}}}
\def\vecA{{\mathbf{A}}}
\def\vecP{{\mathbf{P}}}
\def\vecsigma{{\mathbf{\sigma}}}
\def\hatJn{{\hat{J_\vecn}}}
\def\hatJx{{\hat{J_x}}}
\def\hatJy{{\hat{J_y}}}
\def\hatJz{{\hat{J_z}}}
\def\hatj#1{\hat{J_{#1}}}
\def\hatphi{{\hat{\phi}}}
\def\hatq{{\hat{q}}}
\def\hatpi{{\hat{\pi}}}
\def\vel{\upsilon}
\def\Dint{{\mathcal{D}}}
\def\adag{{\hat{a}^\dagger}}
\def\bdag{{\hat{b}^\dagger}}
\def\cdag{{\hat{c}^\dagger}}
\def\ddag{{\hat{d}^\dagger}}
\def\hata{{\hat{a}}}
\def\hatb{{\hat{b}}}
\def\hatc{{\hat{c}}}
\def\hatd{{\hat{d}}}
\def\hatN{{\hat{N}}}
\def\hatH{{\hat{H}}}
\def\hatp{{\hat{p}}}
\def\Fup{{F^{\mu\nu}}}
\def\Fdown{{F_{\mu\nu}}}
\def\newl{\nonumber \\}
\def\vece{\mathrm{e}}
\def\calM{{\mathcal{M}}}
\def\calT{{\mathcal{T}}}
\def\calR{{\mathcal{R}}}
\def\barpsi{\bar{\psi}}
\def\baru{\bar{u}}
\def\barv{\bar{\upsilon}}
\def\qeq{\stackrel{?}{=}}
\def\torder#1{\mathcal{T}\left(#1\right)}
\def\rorder#1{\mathcal{R}\left(#1\right)}
\def\contr#1#2{\contraction{}{#1}{}{#2}#1#2}
\def\trof#1{\mathrm{Tr}\left(#1\right)}
\def\trace{\mathrm{Tr}}
\def\comm#1{\ \ \ \left(\mathrm{used}\ #1\right)}
\def\tcomm#1{\ \ \ (\text{#1})}
\def\slp{\slashed{p}}
\def\slk{\slashed{k}}
\def\calp{{\mathfrak{p}}}
\def\veccalp{\mathbf{\mathfrak{p}}}
\def\Tthree{T_{\tiny \textcircled{3}}}
\def\pthree{p_{\tiny \textcircled{3}}}
\def\dbar{{\,\mathchar'26\mkern-12mu d}}
\def\erf{\mathrm{erf}}
\def\const{\mathrm{constant}}
\def\pheat{\pfrac p{\ln T}V}
\def\vheat{\pfrac V{\ln T}p}
%%units
\def\fdeg{{^\circ \mathrm{F}}}
\def\cdeg{^\circ \mathrm{C}}
\def\atm{\,\mathrm{atm}}
\def\angstrom{\,\text{\AA}}
\def\SIL{\,\mathrm{L}}
\def\SIkm{\,\mathrm{km}}
\def\SIyr{\,\mathrm{yr}}
\def\SIGyr{\,\mathrm{Gyr}}
\def\SIV{\,\mathrm{V}}
\def\SImV{\,\mathrm{mV}}
\def\SIeV{\,\mathrm{eV}}
\def\SIkeV{\,\mathrm{keV}}
\def\SIMeV{\,\mathrm{MeV}}
\def\SIGeV{\,\mathrm{GeV}}
\def\SIcal{\,\mathrm{cal}}
\def\SIkcal{\,\mathrm{kcal}}
\def\SImol{\,\mathrm{mol}}
\def\SIN{\,\mathrm{N}}
\def\SIHz{\,\mathrm{Hz}}
\def\SIm{\,\mathrm{m}}
\def\SIcm{\,\mathrm{cm}}
\def\SIfm{\,\mathrm{fm}}
\def\SImm{\,\mathrm{mm}}
\def\SInm{\,\mathrm{nm}}
\def\SImum{\,\mathrm{\mu m}}
\def\SIJ{\,\mathrm{J}}
\def\SIW{\,\mathrm{W}}
\def\SIkJ{\,\mathrm{kJ}}
\def\SIs{\,\mathrm{s}}
\def\SIkg{\,\mathrm{kg}}
\def\SIg{\,\mathrm{g}}
\def\SIK{\,\mathrm{K}}
\def\SImmHg{\,\mathrm{mmHg}}
\def\SIPa{\,\mathrm{Pa}}
\def\secpage#1#2{\begin{frame}\bch\bcenter{\bf \Huge #1} \skipline \tbox{\bcenter #2\ecenter}\ecenter\ech\end{frame}}

\newcommand{\field}{\mathscr{F}}
\newcommand{\reals}{\mathbb{R}}
\newcommand{\complexs}{\mathbb{C}}
\newcommand{\ints}{\mathbb{Z}}
%\newcommand{\dim}{\mathrm{dim\ }}
\newcommand{\diag}{\mathrm{diag \ }}
\newcommand{\up}{\uparrow}
\newcommand{\down}{\downarrow}
\newcommand{\su}{\mathfrak{su}}
\newcommand{\so}{\mathfrak{so}}
\newcommand{\tr}{\mathrm{tr\ }}
\newcommand{\card}{\mathrm{card \ }}

\newtheorem{thm}{定理}[section]
\newtheorem{axm}{公理}[section]
\newtheorem{dfn}{定义}[section]

%\cpic{<尺寸>}{<文件名>}}用于生成居中的图片。
\newcommand{\cpic}[2]{
\begin{center}
\includegraphics[scale=#1]{#2}
\end{center}
}

%\cpicn{<尺寸>}{<文件名>}{<注释>}用于生成居中且带有注释的图片，其label为图片名。
\newcommand{\cpicn}[3]
{
\begin{figure}[h!]
\cpic{#1}{#2}
\caption{#3\label{#2}}
\end{figure}
}

\title{RQM Solving Hydrogen Atom}
  \author{Haoting Xu, Zhenjie Liu}
  \date{December 25, 2019}


\begin{document}

\begin{frame}
 
\maketitle
\begin{center}
{\tiny \url{https://github.com/HaotingXu/seminar_lec} }\\
\end{center}
\end{frame}
\section{Runge-Lenz Vector}
\secpage{Runge-Lenz Vector}{$K = \gamma^0 \vec{\Sigma}\cdot \vec{L}+\gamma^0$}
\begin{frame}\frametitle{简要回顾}
  狄拉克方程
  \be
  \left(i\gamma^\mu \partial_\mu-m1_4\right)\psi=0
  \ee
  哈密顿量
  \be
  \hat{H} = i\partial_t = \vec{\alpha}\cdot \vec{p} + \beta m + V(r)
  \ee
  其中
  \be
  \vec{\alpha} =
  \begin{pmatrix}
    &\sigma_i \\
    \sigma_i &
  \end{pmatrix}
  ,\, \beta =
  \begin{pmatrix}
    1&\\
    &-1
  \end{pmatrix}
  \ee
  我们今天要求解$\hat{H} \psi = E\psi $。
  
\end{frame}

\begin{frame}\frametitle{Runge-Lenz Vector}
  我们定义 Runge-Lenz 矢量
  \be
  K = \gamma^0(\Sigma_iL_i + 1)
  \ee
  作为小练习，试证明$[K,H]=0,\,\,[K,J]=0$。
\end{frame}
\begin{frame}\frametitle{对易关系的证明}
  \bea
      [K,H] &=& [\gamma^0\Sigma_i \epsilon_{ijk}x_jP_k,\alpha_m P_m +\gamma^0 m]\\
      &=& [\gamma^0 \Sigma_i\epsilon_{ijk}x_jP_k,\alpha_mP_m] + P_m [\gamma^0,\alpha_m]
  \eea
  计算第一项
  \bea
      \left[\gamma^0 \Sigma_i\epsilon_{ijk}x_jP_k,\alpha_mP_m\right] &=& \gamma^0\alpha_m\Sigma_i\epsilon_{ijk}P_k [x_j,P_m]\\
      &=& i\gamma^0 \Sigma_i \epsilon_{ijk}\alpha_jP_k\\
      &=&i\gamma^0 P_k \epsilon_{ijk}
      \begin{pmatrix}
        & \delta_{ij}+i\epsilon_{ijm}\sigma_{m}\\
        \delta_{ij}+i\epsilon_{ijm}\sigma_{m}&
      \end{pmatrix}\\
      &=&-\gamma^0P_k
      \begin{pmatrix}
        &2\sigma_k\\
        2\sigma_k
      \end{pmatrix}
      \eea
\end{frame}
\begin{frame}\frametitle{对易关系的证明}
     计算第二项
   \be
      P_m[\gamma^0,\alpha_m] = P_k
      \begin{pmatrix}
        0& 2\sigma_k \\
        -2\sigma_k &0
      \end{pmatrix}
   \ee
   故得到$[K,H]=0$。再来考虑$[K,J]$,
   \bea
       [K,J] &=& [\gamma^0(\Sigma_iL_i+1),L_j+\frac{1}{2}\Sigma_j]\\
       &=& \gamma^0\Sigma_i[L_i,L_j] + \frac{1}{2}\gamma^0[\Sigma_i,\Sigma_j]L_i\\
       &=&i\gamma^0 \epsilon_{ijk}(\Sigma_iL_k+\Sigma_kL_i)\\
       &=&0
   \eea
\end{frame}
\begin{frame}\frametitle{量子数}
  假设$K$的本征值为$\kappa$，即$K\psi=\kappa\psi$，则现在描述一个态需要四个量子数
  \be
  |n,j,m_j,\kappa\rangle
  \ee
  普通的量子力学描述氢原子的态
  \be
  |n,l,m_l,m_s\rangle
  \ee
  可见$\kappa$或多或少取代了$l$的地位，现在我们来看看$\kappa$和$l$具体有什么关系。
\end{frame}
\begin{frame}\frametitle{$\kappa$}
  我们来计算$K^2$
  \bea
  K^2 &=&\gamma^0(\Sigma_iL_i + 1)\gamma^0(\Sigma_jL_j + 1)\\
  &=& \Sigma_iL_i \Sigma_jL_j + 2\Sigma_iL_i+1\\
  &=& L_iL_j(\delta_{ij}+i\epsilon_{ijk}\Sigma_k)+2\Sigma_iL_i+1
  \eea
  其中
  \bea
  \epsilon_{ijk}L_iL_j &=& \frac{1}{2}(\epsilon_{ijk}L_iL_j+\epsilon_{jik}L_jL_i)\\
  &=& \frac{1}{2} \epsilon_{ijk}[L_i,L_j]\\
  &=& \frac{i}{2}\epsilon_{ijk}\epsilon_{ijm}L_m\\
  &=& \frac{i}{2}(\delta_{jj}\delta_{km}-\delta_{jm}\delta_{kj})L_m\\
  &=&i L_k
  \eea
\end{frame}
\begin{frame}\frametitle{$\kappa$}
  于是有
  \be
  K^2 = L^2+\Sigma_iL_i+1
  \ee
  对比
  \bea
  J^2 &=& \left( L^2+\Sigma_iL_i +\frac{\Sigma_i\Sigma_i}{4}\right)\\
  &=& L^2+\Sigma_iL_i + \frac{3}{4}
  \eea
  所以有$J^2 = K^2-\frac{1}{4}$，故有$\kappa^2 = j^2+j+\frac{1}{4}$，最终得到
  \be
  \kappa = \pm \left(j+\frac{1}{2}\right)
  \ee
  可见$\kappa$的行为挺像以前的$L^2$的，那么$L^2$发生了什么呢？
\end{frame}
\begin{frame}\frametitle{角动量算符}
  因为$L^2 = J^2-\sigma_iL_i-3/4$，所以要求得$L$需要先求得$\sigma_iL_i$，注意到$K$的定义，有
  \be
  K\psi =
  \begin{pmatrix}
    \sigma_iL_i+1 & 0 \\
     & -\sigma_iL_i-1
  \end{pmatrix}
  \begin{pmatrix}
    \psi_A \\
    \psi_B
  \end{pmatrix}
  = \pm (j+\frac{1}{2})
  \begin{pmatrix}
    \psi_A \\
    \psi_B
  \end{pmatrix}
  \ee
  得到
  \be
  \sigma_iL_i
  \begin{pmatrix}
    \psi_A \\
    \psi_B
  \end{pmatrix}
  =
  \begin{pmatrix}
    \pm(j+\frac{1}{2}\mp 1)\psi_A \\
    \mp(j+\frac{1}{2}\pm 1)\psi_B
  \end{pmatrix}
  \ee
  代入角动量的公式得到
  \be
  L^2
  \begin{pmatrix}
    \psi_A \\
    \psi_B
  \end{pmatrix}
  =
  \begin{pmatrix}
    (j^2+j \pm j\pm \frac{1}{2}+\frac{1}{4})\psi_A \\
    (j^2+j \mp j\mp \frac{1}{2}+\frac{1}{4})\psi_B
  \end{pmatrix}
  \ee
\end{frame}
\begin{frame}\frametitle{角动量算符}
  和以前$l=j\pm \frac{1}{2}$的对比
  \be
  l(l+1) = j^2 + j \pm j\pm \frac{1}{2}+\frac{1}{4}
  \ee
  正好就是上面的结果。但是注意，角动量对于正能量和负能量的态其本征值已经不同，分别记为$l_A,l_B$。这取决于$\kappa$的值，$\kappa$取正，则$l_A = j+\frac{1}{2},\, l_B = j-\frac{1}{2}$，如果$\kappa$取负，则$l_A,l_B$反过来取。
\end{frame}
\section{Raidal Equation}
\secpage{Raidal Equation}{分离变量法}
\begin{frame}\frametitle{求解方程的一般步骤}
  现在我们万事俱备，终于可以来解方程了，我们来回顾一下原子物理中学习到的解方程的一般步骤
  \begin{itemize}
  \item 写出方程。
  \item 写出通解。
  \item 边界条件定系数之间的关系。
  \item 求得能量本征值。
  \end{itemize}
  这里和这个步骤大同小异，我们开始吧。这一小节讲如何写出径向方程。
\end{frame}
\begin{frame}\frametitle{狄拉克方程}
  记$\psi = (\psi_A,\psi_B)^T$，能量本征值方程为
  \be
  (\vec{\alpha}\cdot\vec{p}+\beta m+V)\psi = E\psi
  \ee
  写成矩阵形式
  \be
  \begin{pmatrix}
    0 & \sigma_iP_i \\
    \sigma_iP_i & 0
  \end{pmatrix}
  \begin{pmatrix}
    \psi_A \\
    \psi_B
  \end{pmatrix}
  =
  \begin{pmatrix}
    E-V-m & \\
    & E-V+m
  \end{pmatrix}
  \begin{pmatrix}
    \psi_A \\
    \psi_B
  \end{pmatrix}
  \ee
  这便是要求解的方程，我们首先要化简$\sigma_iP_i$，将它和径向导数扯上关系。以$\hat{x}_i$表示$x_i$方向的单位矢量，有
  \bea
  \sigma_np_n &=& \sigma_i \hat{x}_i\sigma_j \hat{x}_j \sigma_nP_n \\
  &=& \sigma_i \frac{x_i}{r} \sigma_j \frac{x_j}{r} \sigma_n P_n \\
  &=& \frac{1}{r}\frac{\sigma_ix_i}{r}\left( \sigma_j\sigma_nx_jP_n\right) \\
  &=& \frac{1}{r}\frac{\sigma_ix_i}{r}\left( (\delta_{jn}+i\epsilon_{jnk}\sigma_k)x_jP_n\right) \\
  \eea
\end{frame}
\begin{frame}\frametitle{狄拉克方程}
  利用
  \be
  x_jP_j = -ir\frac{\partial}{\partial r}
  \ee
  得到
  \be
  \sigma_np_n = \frac{1}{r}\frac{\sigma_ix_i}{r} \left( -ir \frac{\partial}{\partial r} + i\sigma_iL_i \right)
  \ee
  我们一会将证明，$\frac{\sigma_ix_i}{r},\,i\sigma_iL_i$都和径向部分没什么关系，故我们设
  \be
  \psi_A = g(r) \mathcal{Y}^{m_j}_{jl_A},\,\,\, \psi_B = if(r)\mathcal{Y}^{m_j}_{jl_B}
  \ee
  我们刚刚说过，$\psi_A$和$\psi_B$都是$L^2$的本征态(但是他们拼起来不是)，我们之前学过角动量的本征函数为$Y_{lm}$,故有
  \be
  \mathcal{Y}^{m_j}_{jl_A} = \alpha Y_{l_A,m_j-1/2}\chi_++\beta Y_{l_A,m_j+1/2}\chi_-
  \ee
\end{frame}
\begin{frame}\frametitle{狄拉克方程}
  将上面的分离变量带入矩阵形式的狄拉克方程中，得到
  \bea
  &\frac{1}{r}\frac{\sigma_ix_i}{r} \left( -ir \frac{\partial}{\partial r} + i\sigma_iL_i \right)
  \begin{pmatrix}
    if(r)\mathcal{Y}^{m_j}_{jl_B} \\
    g(r) \mathcal{Y}^{m_j}_{jl_A}
  \end{pmatrix}
  \\&=
  \begin{pmatrix}
    E-V-m & \\
    &E-V+m
  \end{pmatrix}
  \begin{pmatrix}
    g(r)\mathcal{Y}^{m_j}_{jl_A} \\
    if(r)\mathcal{Y}^{m_j}_{jl_B}
  \end{pmatrix}
  \eea
  我们接下来研究算符$\frac{\sigma_ix_i}{r},\,i\sigma_iL_i$的作用，我们将会惊奇的发现，它们让角度部分消掉了。
\end{frame}
\begin{frame}\frametitle{$\sigma_iL_i$}
  首先因为$\sigma_iL_i$与$K$直接相关，所以因为$K\psi = \kappa \psi$，有
  \be
  \begin{pmatrix}
    \sigma_iL_i + 1 & 0 \\
    0& -\sigma_iL_i-1
  \end{pmatrix}
  \begin{pmatrix}
    \psi_A \\
    \psi_B
  \end{pmatrix}
  =\kappa
  \begin{pmatrix}
    \psi_A \\
    \psi_B
  \end{pmatrix}
  \ee
  得到
  \bea
  \sigma_iL_i \psi_B = (\kappa-1)\psi_B\\
  \sigma_iL_i \psi_A = (-\kappa-1) \psi_A
  \eea
\end{frame}
\begin{frame}\frametitle{$\frac{\sigma_ix_i}{r}$}
  首先，请尝试证明$[\sigma_ix_i,J_i]=0,\,[\frac{\sigma_ix_i}{r},H]=0$，即这个算符不改变$j$和$m_j$和$n$。又因为它是一个pseudoscalar，所以在坐标变换$\hat{\pi}$下，有
  \be
  \hat{\pi}\left(\frac{\sigma_ix_i}{r} \psi\right) \sim (-1)\frac{\sigma_ix_i}{r}\hat{\pi}\psi \sim (-1)^{l+1}\psi
  \ee
  可见，这个算符改变了态的宇称。我们知道态的宇称由量子数$l$来决定，所以说，我们有
  \be
  \frac{\sigma_ix_i}{r} \mathcal{Y}^{m_j}_{jl_A} = C \mathcal{Y}^{m_j}_{jl_B}
  \ee
  又因为
  \be
  \left(\frac{\sigma_ix_i}{r}\right)^2=1
  \ee
  所以$C=e^{i\delta}$，为了方便，我们取$C=-1$，故有
  \be
  \frac{\sigma_ix_i}{r} \mathcal{Y}^{m_j}_{jl_A} = - \mathcal{Y}^{m_j}_{jl_B}
  \ee
\end{frame}
\begin{frame}\frametitle{Radial Equation}
  将两个算符带入狄拉克方程疯狂化简，我们发现角度部分竟然约掉了，只剩下径向部分
  \bea
  \left(-\partial_r + \frac{(-\kappa-1)}{r}\right) f(r) &=& (E-V-m)g(r)\\
  \left(-\partial_r + \frac{(\kappa-1)}{r}\right) g(r) &=& -(E-V+m)f(r)
  \eea
  凭借着物理学家的直觉，定义$F=rf,\,G=rg$，带入方程中，得到径向方程
  \bea
  \frac{\partial F}{\partial r}+ \frac{\kappa F}{r}+ (E-V-m)G&=&0\\
  \frac{\partial G}{\partial r}- \frac{\kappa G}{r}- (E-V+m)F&=&0
  \eea
  注意，上式对于任何的$V(r)$都成立。
\end{frame}
\section{Energy}
\begin{frame}\frametitle{变量替换}
  现在引入库仑势
  \be
  V = -\frac{Z\alpha}{r}\hbar c = -\frac{Z\alpha}{r}
  \ee
  我们定义$k_1= m+E,\,k^2 = m-E,\rho = \sqrt{k_1k_2}r$，带入方程中得到
  \bea
  \left(\frac{\partial}{\partial \rho}- \frac{\kappa}{\rho}\right)F - \left(\sqrt{\frac{k_2}{k_1}}-\frac{Z\alpha}{\rho}\right) G &=& 0 \\
  \left(\frac{\partial}{\partial \rho}+ \frac{\kappa}{\rho}\right)G - \left(\sqrt{\frac{k_2}{k_1}}+\frac{Z\alpha}{\rho}\right) F &=& 0
  \eea
  先来研究$\rho\rightarrow\infty$时方程的行为，当$\rho\rightarrow \infty$时，有
  \be
  \frac{\partial^2 F}{\partial \rho^2}-F =0
  \ee
\end{frame}
\begin{frame}\frametitle{级数解}
  所以我们假设
  \be
  F = \rho^se^{-\rho}\sum_{m=0}^{\infty} a_m\rho^m ,\,\,\, G = \rho^se^{-\rho}\sum_{m=0}^{\infty} b_m\rho^m
  \ee
  带入方程中并化简得到
  \bea
  -a_m + (s+m+1-\kappa)a_{m+1} - \sqrt{\frac{k_2}{k_1}}b_m +Z\alpha b_{m+1}&=&0\\
  -b_m + (s+m+1+\kappa)b_{m+1} - \sqrt{\frac{k_1}{k_2}}a_m -Z\alpha a_{m+1}&=&0
  \eea
  取$m=-1$，并要求级数没有负次方项
  \bea
  (s-\kappa)a_0 + Z\alpha b_0 &=&0 \\
  (s+\kappa)b_0 -Z\alpha a_0 &=&0
  \eea
\end{frame}

\begin{frame}\frametitle{求出$s$}
  若使得上面方程组有解，则有
  \be
  s = \pm\sqrt{\kappa^2 - Z^2\alpha^2}
  \ee
  因为$Z\alpha$很小，$\kappa$是一个大于1的整数，所以上式挺合理的。
\end{frame}
\begin{frame}\frametitle{令级数终止}
  现在我们让级数终止于$m=N$，即$a_{N+1}=b_{N+1}=0$，带入递推关系有
  \be
  a_N = -\sqrt{\frac{k_2}{k_1}}b_N,\,\,\, b_N = -\sqrt{\frac{k_1}{k_2}}a_N
  \ee
  现在令$m=N-1$，就可以得到能量本征值了。
  \bea
  -a_{N-1}+(s+N-\kappa)a_N - \sqrt{\frac{k_2}{k_1}}b_{N-1}+Z\alpha b_N &=&0\\
  -b_{N-1}+(s+N+\kappa)a_N - \sqrt{\frac{k_1}{k_2}}a_{N-1}+Z\alpha a_N &=&0
  \eea
  聪明的你一定发现了这两个方程是一个方程。带入上面的关系，消掉$a_N$，得到
  \be
  \sqrt{\frac{k_1}{k_2}}-\sqrt{\frac{k_2}{k_1}} = \frac{2(s+N)}{Z\alpha}
  \ee
\end{frame}
\begin{frame}\frametitle{能量本征值}
  将有关定义带入，得到
  \bea
  E &=& \frac{m}{\sqrt{1+\frac{Z^2\alpha^2}{(N+s)^2}}} \\
  &=& \frac{mc^2}{\sqrt{1+\frac{Z^2\alpha^2}{(N+\sqrt{\kappa^2-Z^2\alpha^2})^2}}}
  \eea
\end{frame}
\section{Expasion}
\begin{frame}\frametitle{展开}
  因为$Z\alpha$一般特别小，所以上式可按照$Z\alpha$展开，代入mathematica中得到
  \be
  E \simeq mc^2 - \frac{mc^2(Z\alpha)^2}{(N+j+1/2)^2}-\frac{1}{2}\frac{(Z\alpha)^4mc^2}{(N+j+1/2)^4}\left(\frac{N+j+1/2}{j+1/2}-\frac{3}{4}\right)+\cdots
  \ee
  从上式我们可以认出来原来的主量子数$n=N+j+1/2$。而四次方项正是我们之前由微扰论算出的精细结构。
\end{frame}
\begin{frame}\frametitle{数值结果}
  首先对于氢原子，计算一下数值如下图。可见结果差别很小，这是由于氢原子中电子运动的能量为$10\mathrm{eV}$，而电子质量大概$0.511\mathrm{MeV}$，所以速度并不是很大。
  \cpic{0.3}{hydrogen}
\end{frame}
\begin{frame}\frametitle{$Z$很大}
  下表给出了$Z=100$的类氢原子的结果，可见差异很大。
  \cpic{0.3}{100_hydro}
\end{frame}
\section{goodbye}
\begin{frame}\frametitle{Merry Christmas}
  不知道说什么好了，Merry Christmas!
  \cpic{0.3}{christmas}
\end{frame}




  

  
\end{document}
